\scanpage{49}
% archive foliation: 48
$U$ sss $x\in U^\nu$, \deleted{donc} \unsure{or implique|or ceci implique} $x\in U_i^\nu$ pour $i$ convenable i.e. les $x_\alpha\in U_i$. \unsure{O.K.}

On est ainsi ramené au cas où $X$ est \uline{affine} sur $Y$. Mais alors, on est $=$ ramené au cas où $X=\unsure{\operatorname{Spec}|\operatorname{Sym}}(\unsure{\uline F})$, où $F$ est un faisceau loc. libre de rang $n$ sur $T$. On sait qu'alors
\[
f_*(X/T)=\unsure{\operatorname{Spec}|\operatorname{Sym}}(f_*\unsure{\check F|\uline F}).
\]

\deleted{\uline{\unsure{Corollaire}} $\unsure{f_*}$ \unsure{de}}

\uline{Remarque} Il n'est pas vrai en général que la dim. des fibres de $\unsure{f_*|f^*}(X/T)/S$ soit majorée par $\unsure{\nu|d}$ fois la dim des fibres de \deleted{$\unsure{f_*}$} $X/T$. Ex:
$S=\operatorname{Spec}(k)$, $T$\deleted{$=\operatorname{Spec}k[\varepsilon]/\unsure{\varepsilon^2}$} est purement radiciel sur $S$, $X=T\times_S T$ i.e. $f_*(X/T)=\uline{\operatorname{Hom}}_S(T,T)$. Ce dernier peut être de \uline{dim} \inserted{$>0$} \deleted{$\unsure{d}$} sur $k$. Exemple: le groupe des automorphismes de la $k$-algèbre $\unsure{A|k}$ est de dim $>0$.
\margin{\unsure{voir précisé|voir préciser} plus bas}
